\section{Testes de continuidade via curvas}\index{teste das curvas}\index{curvas!teste de continuidade}

Funções de várias variáveis podem ser difíceis de visualizar diretamente.
Uma técnica muito útil é aproximar o ponto por curvas.
Ao compor uma função de várias variáveis com uma curva, obtemos uma função de uma variável, que costuma ser mais simples de estudar.

O princípio básico é o seguinte: se uma função possui limite em um ponto, então esse mesmo limite aparece ao longo de qualquer curva que se aproxima desse ponto sem ficar presa nele.

\begin{proposition}\index{limite!ao longo de curvas}
    Sejam:
    \begin{itemize}
        \item $A\subseteq\mathbb R^n$ um conjunto,
        \item $f:A\to\mathbb R^m$ uma função,
        \item $p\in\mathbb R^n$ um ponto de acumulação de $A$,
        \item $L\in\mathbb R^m$,
        \item $I\subseteq\mathbb R$ um conjunto,
        \item $t_0\in\mathbb R$ um ponto de acumulação de $I$, e
        \item $\gamma:I\to A$ uma função.
    \end{itemize}

    Suponha que:
    \begin{itemize}
        \item $\lim_{t\to t_0}\gamma(t)=p$;
        \item existe $\eta>0$ tal que, para todo $t\in I\setminus\{t_0\}$,
        \[
            |t-t_0|<\eta
            \quad\Longrightarrow\quad
            \gamma(t)\neq p;
        \]
        \item $\lim_{x\to p}f(x)=L$.
    \end{itemize}

    Então
    \[
        \lim_{t\to t_0}f(\gamma(t))=L.
    \]
\end{proposition}

\begin{proof}
    Fixe $\epsilon>0$.

    Como $\lim_{x\to p}f(x)=L$, existe $\delta_f>0$ tal que, para todo $x\in A\setminus\{p\}$,
    \[
        \text{se }d(x,p)<\delta_f,\text{ então }d(f(x),L)<\epsilon.
    \]

    Como $\lim_{t\to t_0}\gamma(t)=p$, existe $\delta_\gamma>0$ tal que, para todo $t\in I\setminus\{t_0\}$,
    \[
        \text{se }|t-t_0|<\delta_\gamma,\text{ então }d(\gamma(t),p)<\delta_f.
    \]

    Tome
    \[
        \delta=\min\{\eta,\delta_\gamma\}.
    \]

    Se $t\in I\setminus\{t_0\}$ e $|t-t_0|<\delta$, então $\gamma(t)\neq p$ e $d(\gamma(t),p)<\delta_f$.
    Como $\gamma(t)\in A$, temos $\gamma(t)\in A\setminus\{p\}$.
    Portanto,
    \[
        d(f(\gamma(t)),L)<\epsilon.
    \]

    Logo,
    \[
        \lim_{t\to t_0}f(\gamma(t))=L.
    \]
\end{proof}

\begin{corollary}[Teste negativo por curvas]\index{limite!teste negativo por curvas}
    Seja $f:A\to\mathbb R^m$, com $A\subseteq\mathbb R^n$, e seja $p$ um ponto de acumulação de $A$.

    Se existir uma curva $\gamma:I\to A$ tal que $\gamma(t)\to p$ quando $t\to t_0$, $\gamma(t)\neq p$ para $t\neq t_0$ suficientemente próximo de $t_0$, e $f\circ\gamma$ não possuir limite em $t_0$, então $f$ não possui limite em $p$.

    Além disso, se existirem duas curvas $\gamma_1:I_1\to A$ e $\gamma_2:I_2\to A$ que se aproximam de $p$ e tais que
    \[
        \lim_{t\to t_1}f(\gamma_1(t))
        \neq
        \lim_{t\to t_2}f(\gamma_2(t)),
    \]
    então $f$ não possui limite em $p$.
\end{corollary}

\begin{proof}
    Se $f$ possuísse limite em $p$, então, pela proposição anterior, toda curva que se aproxima de $p$ satisfazendo as hipóteses teria, após composição com $f$, esse mesmo limite.
    Portanto, uma curva sem limite, ou duas curvas com limites diferentes, impossibilitam a existência do limite de $f$ em $p$.
\end{proof}

\begin{corollary}\index{continuidade!teste negativo por curvas}
    Seja $A\subseteq\mathbb R^n$, seja $f:A\to\mathbb R^m$ uma função e seja $p\in A$.

    Se $f$ é contínua em $p$, então, para toda curva contínua $\gamma:I\to A$ e todo $t_0\in I$ tal que $\gamma(t_0)=p$, a função
    \[
        f\circ\gamma:I\to\mathbb R^m
    \]
    é contínua em $t_0$.

    Consequentemente, se existe uma curva contínua $\gamma:I\to A$ com $\gamma(t_0)=p$ tal que $f\circ\gamma$ não é contínua em $t_0$, então $f$ não é contínua em $p$.
\end{corollary}

\begin{proof}
    Fixe $\epsilon>0$.
    Como $f$ é contínua em $p$, existe $\delta_f>0$ tal que, para todo $x\in A$,
    \[
        \text{se }d(x,p)<\delta_f,\text{ então }d(f(x),f(p))<\epsilon.
    \]

    Como $\gamma$ é contínua em $t_0$ e $\gamma(t_0)=p$, existe $\delta_\gamma>0$ tal que, para todo $t\in I$,
    \[
        \text{se }|t-t_0|<\delta_\gamma,\text{ então }d(\gamma(t),p)<\delta_f.
    \]

    Assim, para todo $t\in I$ com $|t-t_0|<\delta_\gamma$,
    \[
        d(f(\gamma(t)),f(\gamma(t_0)))=d(f(\gamma(t)),f(p))<\epsilon.
    \]
    Portanto, $f\circ\gamma$ é contínua em $t_0$.

    A última afirmação é a contrapositiva da primeira.
\end{proof}

\begin{example}
    Seja $f:\mathbb R^2\setminus\{(0,0)\}\to\mathbb R$ dada por
    \[
        f(x,y)=\frac{xy}{x^2+y^2}.
    \]

    Vamos mostrar que $f$ não possui limite em $(0,0)$.

    Considere as curvas
    \[
        \gamma_1(t)=(t,0)
        \qquad\text{e}\qquad
        \gamma_2(t)=(t,t),
    \]
    definidas para $t\neq 0$.
    Ambas se aproximam de $(0,0)$ quando $t\to 0$.

    Ao longo de $\gamma_1$, temos
    \[
        f(\gamma_1(t))=f(t,0)=0.
    \]
    Logo,
    \[
        \lim_{t\to 0}f(\gamma_1(t))=0.
    \]

    Ao longo de $\gamma_2$, temos
    \[
        f(\gamma_2(t))=f(t,t)=\frac{t^2}{2t^2}=\frac{1}{2}.
    \]
    Logo,
    \[
        \lim_{t\to 0}f(\gamma_2(t))=\frac{1}{2}.
    \]

    Como obtivemos limites diferentes ao longo de duas curvas que se aproximam de $(0,0)$, concluímos que
    \[
        \lim_{(x,y)\to(0,0)}\frac{xy}{x^2+y^2}
    \]
    não existe.
\end{example}

\begin{example}
    Verificar apenas retas pode ser insuficiente.

    Defina $F:\mathbb R^2\to\mathbb R$ por
    \[
        F(x,y)=
        \begin{cases}
            \dfrac{x^2y}{x^4+y^2}, & \text{se }(x,y)\neq(0,0),\\
            0, & \text{se }(x,y)=(0,0).
        \end{cases}
    \]

    Ao longo de qualquer reta passando pela origem, obtemos limite $0$.
    De fato, se a reta é $y=ax$, então, para $a\neq 0$,
    \[
        F(t,at)
        =
        \frac{x^2y}{x^4+y^2}\bigg|_{(x,y)=(t,at)}
        =
        \frac{at^3}{t^4+a^2t^2}
        =
        \frac{at}{t^2+a^2}
        \longrightarrow 0.
    \]
    Se $a=0$, então $F(t,0)=0$.
    A reta vertical $x=0$ também dá $F(0,t)=0$.

    No entanto, ao longo da parábola $\gamma(t)=(t,t^2)$,
    \[
        F(\gamma(t))
        =
        F(t,t^2)
        =
        \frac{t^2t^2}{t^4+t^4}
        =
        \frac{1}{2}
    \]
    para todo $t\neq 0$.

    Portanto, $F$ não é contínua em $(0,0)$.
    Esse exemplo mostra que testar somente retas pode sugerir um comportamento falso.
\end{example}

O teste negativo por curvas é frequentemente o mais usado: encontramos uma curva ruim e concluímos que o limite não existe ou que a função não é contínua.
Mas existe também uma versão completa do teste, desde que aceitemos caminhos de teste cujos domínios sejam subconjuntos de $\mathbb R$, e não apenas intervalos.

\begin{definition}\index{caminho de teste}
    Seja $A\subseteq\mathbb R^n$ e seja $p\in\mathbb R^n$.

    Um \emph{caminho de teste em $A$ chegando a $p$} é uma função $\gamma:D\to A$, em que $D\subseteq\mathbb R$ possui $0$ como ponto de acumulação e
    \[
        \lim_{t\to 0}\gamma(t)=p.
    \]

    Dizemos que o caminho de teste é \emph{perfurado} se $\gamma(t)\neq p$ para todo $t\in D$.
\end{definition}

\begin{proposition}[Teste completo por caminhos]\index{limite!teste completo por caminhos}
    Seja $A\subseteq\mathbb R^n$, seja $f:A\to\mathbb R^m$, seja $p$ um ponto de acumulação de $A$ e seja $L\in\mathbb R^m$.

    Então $\lim_{x\to p}f(x)=L$ se, e somente se, para todo caminho de teste perfurado $\gamma:D\to A$ chegando a $p$, temos
    \[
        \lim_{t\to 0}f(\gamma(t))=L.
    \]
\end{proposition}

\begin{proof}
    A ida é a mesma ideia da proposição sobre limites ao longo de curvas.
    Se $\lim_{x\to p}f(x)=L$ e $\gamma:D\to A$ é um caminho de teste perfurado chegando a $p$, então, dado $\epsilon>0$, existe $\delta_f>0$ tal que
    \[
        x\in A\setminus\{p\},\ d(x,p)<\delta_f
        \quad\Longrightarrow\quad
        d(f(x),L)<\epsilon.
    \]
    Como $\gamma(t)\to p$, existe $\delta_\gamma>0$ tal que
    \[
        t\in D,\ 0<|t|<\delta_\gamma
        \quad\Longrightarrow\quad
        d(\gamma(t),p)<\delta_f.
    \]
    Como $\gamma$ é perfurado, $\gamma(t)\in A\setminus\{p\}$ para todo $t\in D$.
    Logo,
    \[
        d(f(\gamma(t)),L)<\epsilon
    \]
    sempre que $t\in D$ e $0<|t|<\delta_\gamma$.
    Portanto,
    \[
        \lim_{t\to 0}f(\gamma(t))=L.
    \]

    Provemos a recíproca por contraposição.
    Suponha que $L$ não é limite de $f$ em $p$.
    Então existe $\epsilon_0>0$ tal que, para todo $\delta>0$, podemos encontrar $x\in A\setminus\{p\}$ satisfazendo
    \[
        d(x,p)<\delta
        \qquad\text{e}\qquad
        d(f(x),L)\geq\epsilon_0.
    \]

    Para cada inteiro $k\geq 1$, aplique essa afirmação com $\delta=\frac{1}{k}$.
    Obtemos $x_k\in A\setminus\{p\}$ tal que
    \[
        d(x_k,p)<\frac{1}{k}
        \qquad\text{e}\qquad
        d(f(x_k),L)\geq\epsilon_0.
    \]

    Defina
    \[
        D=\left\{\frac{1}{k}:k\in\mathbb N,\ k\geq 1\right\}
    \]
    e $\gamma:D\to A$ por
    \[
        \gamma\left(\frac{1}{k}\right)=x_k.
    \]

    O conjunto $D$ possui $0$ como ponto de acumulação, e $\gamma(t)\to p$ quando $t\to 0$ em $D$, pois $d(x_k,p)<\frac{1}{k}$.
    Além disso, $\gamma$ é perfurado, pois $x_k\neq p$ para todo $k$.

    Porém,
    \[
        d(f(\gamma(1/k)),L)=d(f(x_k),L)\geq\epsilon_0
    \]
    para todo $k$.
    Portanto, $f\circ\gamma$ não tende a $L$ quando $t\to 0$ em $D$.

    Assim, a hipótese sobre todos os caminhos de teste perfurados implica necessariamente que
    \[
        \lim_{x\to p}f(x)=L.
    \]
\end{proof}

\begin{corollary}[Teste completo de continuidade]\index{continuidade!teste completo por caminhos}
    Seja $A\subseteq\mathbb R^n$, seja $f:A\to\mathbb R^m$ e seja $p\in A$.

    Então $f$ é contínua em $p$ se, e somente se, para todo conjunto $D\subseteq\mathbb R$ com $0\in D$ e $0$ ponto de acumulação de $D$, e para toda função $\gamma:D\to A$ contínua em $0$ tal que $\gamma(0)=p$, a função $f\circ\gamma:D\to\mathbb R^m$ é contínua em $0$.
\end{corollary}

\begin{proof}
    Se $p$ é ponto isolado de $A$, então $f$ é contínua em $p$.
    Além disso, se $\gamma:D\to A$ é contínua em $0$ e $\gamma(0)=p$, então, para $t$ suficientemente próximo de $0$, temos $\gamma(t)=p$.
    Logo, $f\circ\gamma$ é contínua em $0$.

    Suponha, então, que $p$ é ponto de acumulação de $A$.

    Se $f$ é contínua em $p$, então $f\circ\gamma$ é contínua em $0$ pela continuidade de composições.
    De fato, dado $\epsilon>0$, existe $\delta_f>0$ tal que
    \[
        x\in A,\ d(x,p)<\delta_f
        \quad\Longrightarrow\quad
        d(f(x),f(p))<\epsilon.
    \]
    Como $\gamma$ é contínua em $0$ e $\gamma(0)=p$, existe $\delta_\gamma>0$ tal que
    \[
        t\in D,\ |t|<\delta_\gamma
        \quad\Longrightarrow\quad
        d(\gamma(t),p)<\delta_f.
    \]
    Logo, para $t\in D$ com $|t|<\delta_\gamma$,
    \[
        d(f(\gamma(t)),f(\gamma(0)))=d(f(\gamma(t)),f(p))<\epsilon.
    \]
    Portanto, $f\circ\gamma$ é contínua em $0$.

    Reciprocamente, suponha que $f$ não é contínua em $p$.
    Como $p$ é ponto de acumulação de $A$, temos que $f(p)$ não é limite de $f$ em $p$.
    Pela prova da proposição anterior, existe um caminho de teste perfurado $\alpha:D_0\to A$ chegando a $p$ tal que $f(\alpha(t))$ não tende a $f(p)$ quando $t\to 0$ em $D_0$.

    Defina $D=D_0\cup\{0\}$ e $\gamma:D\to A$ por
    \[
        \gamma(0)=p
        \qquad\text{e}\qquad
        \gamma(t)=\alpha(t)\text{ se }t\in D_0.
    \]
    Então $\gamma$ é contínua em $0$ e $\gamma(0)=p$.
    Porém $f\circ\gamma$ não é contínua em $0$, pois não se aproxima de $f(p)=(f\circ\gamma)(0)$ ao longo dos pontos de $D_0$.

    Portanto, se $f\circ\gamma$ é contínua em $0$ para todo caminho de teste contínuo $\gamma$ com $\gamma(0)=p$, então $f$ é contínua em $p$.
\end{proof}

\section*{Exercícios}

\begin{exercise}
    Use curvas para mostrar que o limite abaixo não existe:
    \[
        \lim_{(x,y)\to(0,0)}\frac{x^2-y^2}{x^2+y^2}.
    \]
\end{exercise}

\begin{exercise}
    Seja $f:\mathbb R^2\to\mathbb R$ dada por
    \[
        f(x,y)=
        \begin{cases}
            \dfrac{y^2}{x^2+y^2}, & \text{se }(x,y)\neq(0,0),\\
            0, & \text{se }(x,y)=(0,0).
        \end{cases}
    \]
    Use curvas para decidir se $f$ é contínua em $(0,0)$.
\end{exercise}

\begin{exercise}
    Seja $g:\mathbb R^2\to\mathbb R$ dada por
    \[
        g(x,y)=
        \begin{cases}
            \dfrac{x^2y}{x^4+y^2}, & \text{se }(x,y)\neq(0,0),\\
            0, & \text{se }(x,y)=(0,0).
        \end{cases}
    \]
    Verifique que, ao longo de toda reta passando pela origem, o limite de $g$ é $0$.
    Em seguida, encontre uma curva que mostre que $g$ não é contínua em $(0,0)$.
\end{exercise}

\begin{exercise}
    Seja $F:\mathbb R^2\to\mathbb R^2$ dada por
    \[
        F(x,y)=
        \begin{cases}
            \left(\dfrac{xy}{x^2+y^2}, x+y\right), & \text{se }(x,y)\neq(0,0),\\
            (0,0), & \text{se }(x,y)=(0,0).
        \end{cases}
    \]
    Use curvas e o critério por coordenadas para mostrar que $F$ não é contínua em $(0,0)$.
\end{exercise}

\begin{exercise}
    Seja $f:A\to\mathbb R^m$ uma função contínua em $p\in A$, e seja $\gamma:I\to A$ uma curva contínua tal que $\gamma(t_0)=p$.
    Mostre diretamente pela definição de continuidade que $f\circ\gamma$ é contínua em $t_0$.
\end{exercise}
